%!TEX root = ../../thesis.tex
\section{関連研究}
\label{sec:related}

% TODO: 視覚と行動のend-to-end学習（条件付き模倣学習など）と，トポロジカルマップを用いた視覚ナビゲーションの研究を書く
%       （春山の修論 1.2 節で挙げている文献を参照）．

\subsection{Vision-and-Language Navigation}
\label{sec:related_vln}

関連研究として，言語で記述された指示に従って移動するVision-and-Language Navigation（VLN）が挙げられる．
LM-Navは，大規模言語モデルで指示から目印を抽出し，視覚言語モデルで目印を事前に構築したトポロジカルグラフ上の画像と対応付け，視覚ナビゲーションモデルで移動する\cite{lmnav}．
NaVILAは，視覚言語モデルが「前に75\,cm進む」のような中間的な言語の行動を出力し，それを脚ロボットの移動方策が実行する\cite{navila}．
LM-Nav\cite{lmnav}やNaVILA\cite{navila}のように，言語による指示で未学習の環境を移動すること自体はすでに実現されつつある．
しかし，これらの手法は数十億以上のパラメータを持つ言語モデルや視覚言語モデルを用いており，それらの推論をロボット外部の計算機で行っている\cite{lmnav,navila}．

\subsection{本研究の位置付け}
\label{sec:related_position}

本研究では，シナリオの解釈をルールベースで処理し，深層学習モデルを通路分類と経路追従に限定することで，用いるモデルのパラメータ数を合計約7500万に抑え，ロボットに搭載可能な計算機での動作を想定したシステムを構成する．
また，実行時に用いるセンサは単眼RGBカメラのみである．
視覚ナビゲーション基盤モデルと比べると，ゴール画像の代わりにシナリオから決まる目標方向を与える点が異なり，走行対象の環境を事前に走行する必要がない．
VLNと比べると，自由な言語の指示は扱えないが，定型的なシナリオに限定することで，大規模な言語モデルを用いずに経路を指示できる．
